\sectionheading{5.32 Variable ও Data Type}

Program-এ বিভিন্ন ধরনের Data নিয়ে কাজ করতে হয়।

কোনো Data বা Value-কে একটি নাম দিয়ে Program-এ ব্যবহার করার
জন্য \textbf{Variable} এবং Dataটি কী ধরনের তা বোঝানোর জন্য
\textbf{Data Type} ব্যবহৃত হয়।

%================================================
% Variable Definition
%================================================

\booksubheading{Variable কী?}

\begin{definitionbox}

\textbf{\color{darkTeal}সংজ্ঞা}

Program-এ কোনো Data বা Value-কে একটি নামের মাধ্যমে
সংরক্ষণ ও ব্যবহার করার ব্যবস্থাকে সাধারণভাবে
\textbf{Variable} বলা হয়।

\end{definitionbox}

উদাহরণ—

\begin{codebox}

\begin{english}
\ttfamily
age = 18
\end{english}

\end{codebox}

এখানে—

\begin{tabularx}
{\linewidth}
{
|>{\RaggedRight\arraybackslash}p{0.30\linewidth}
|Y|
}

\hline

\rowcolor{softBlue}

\multicolumn{1}{|c|}{\textbf{অংশ}}
&
\multicolumn{1}{c|}{\textbf{অর্থ}}

\\

\hline

\textenglish{\ttfamily age}
&
Variable-এর নাম
\\

\hline

\textenglish{\ttfamily =}
&
Assignment Operator
\\

\hline

\textenglish{\ttfamily 18}
&
Variable-এর সঙ্গে যুক্ত Value
\\

\hline

\end{tabularx}

\begin{rememberbox}

\textbf{\color{deepOrange}সহজভাবে মনে রাখবে}

\begin{center}

\textbf{Variable}
\(\longrightarrow\)
\textbf{Data/Value-এর নামযুক্ত ধারক}

\end{center}

\end{rememberbox}

%================================================
% Why Variable
%================================================

\booksubheading{Variable কেন প্রয়োজন?}

ধরা যাক—

\begin{codebox}

\begin{english}
\ttfamily
price = 500
\end{english}

\end{codebox}

এখন Program-এর বিভিন্ন জায়গায় সরাসরি
\textenglish{\ttfamily 500}
না লিখে
\textenglish{\ttfamily price}
নামটি ব্যবহার করা যায়।

উদাহরণ—

\begin{codebox}

\begin{english}
\ttfamily
print(price)
\end{english}

\end{codebox}

\textbf{Output:}

\begin{codebox}

\begin{english}
\ttfamily
500
\end{english}

\end{codebox}

এতে Program লেখা, বোঝা ও পরিবর্তন করা সহজ হয়।

%================================================
% Data Type Definition
%================================================

\booksubheading{Data Type কী?}

\begin{definitionbox}

\textbf{\color{darkTeal}সংজ্ঞা}

কোনো Data বা Value কী ধরনের এবং তার উপর কী ধরনের
Operation করা যাবে—তা নির্দেশকারী শ্রেণিকে
\textbf{Data Type} বলা হয়।

\end{definitionbox}

উদাহরণ—

\begin{codebox}

\begin{english}
\ttfamily
age = 18\\
name = "Rahim"
\end{english}

\end{codebox}

এখানে—

\begin{itemize}

\item
\textenglish{\ttfamily 18}
একটি পূর্ণসংখ্যা।

\item
\textenglish{\ttfamily "Rahim"}
একটি Text/String ধরনের Data।

\end{itemize}

%================================================
% Why Data Type
%================================================

\booksubheading{Data Type কেন প্রয়োজন?}

Computer-এর কাছে নিচের Valueগুলো একই ধরনের Data নয়—

\begin{codebox}

\begin{english}
\ttfamily
20\\
20.5\\
"20"\\
True
\end{english}

\end{codebox}

\begin{tabularx}
{\linewidth}
{
|>{\RaggedRight\arraybackslash}p{0.30\linewidth}
|Y|
}

\hline

\rowcolor{softBlue}

\multicolumn{1}{|c|}{\textbf{Value}}
&
\multicolumn{1}{c|}{\textbf{সাধারণ ধারণা}}

\\

\hline

\textenglish{\ttfamily 20}
&
পূর্ণসংখ্যা
\\

\hline

\textenglish{\ttfamily 20.5}
&
দশমিক সংখ্যা
\\

\hline

\textenglish{\ttfamily "20"}
&
Text/String
\\

\hline

\textenglish{\ttfamily True}
&
Logical/Boolean Value
\\

\hline

\end{tabularx}

তাই Program-এ Data-এর ধরন জানা গুরুত্বপূর্ণ।

%================================================
% Relation
%================================================

\booksubheading{Variable ও Data Type-এর সম্পর্ক}

Variable কোনো Value-কে নাম দিয়ে ব্যবহার করতে সাহায্য করে এবং
Data Type সেই Value-এর ধরন নির্দেশ করে।

\begin{examplebox}{উদাহরণ ১}

\begin{codebox}

\begin{english}
\ttfamily
age = 18
\end{english}

\end{codebox}

এখানে—

\begin{center}

\textbf{Variable}
\(\longrightarrow\)
\textenglish{\ttfamily age}
\qquad
\textbf{Value}
\(\longrightarrow\)
\textenglish{\ttfamily 18}
\qquad
\textbf{Data Type}
\(\longrightarrow\)
Integer

\end{center}

\end{examplebox}

\begin{examplebox}{উদাহরণ ২}

\begin{codebox}

\begin{english}
\ttfamily
name = "Karim"
\end{english}

\end{codebox}

এখানে—

\begin{center}

\textbf{Variable}
\(\longrightarrow\)
\textenglish{\ttfamily name}
\qquad
\textbf{Value}
\(\longrightarrow\)
\textenglish{\ttfamily "Karim"}
\qquad
\textbf{Data Type}
\(\longrightarrow\)
String

\end{center}

\end{examplebox}

\begin{rememberbox}

\textbf{\color{deepOrange}দ্রুত সম্পর্ক}

\begin{center}

\textbf{Variable Name}
\(\longrightarrow\)
\textbf{Value}
\(\longrightarrow\)
\textbf{Data Type}

\end{center}

\end{rememberbox}

%================================================
% Different Types of Data
%================================================

\booksubheading{একই Program-এ বিভিন্ন ধরনের Data}

নিচের Programটি লক্ষ্য করো—

\begin{codebox}

\begin{english}
\ttfamily
name = "Rahim"\\
age = 18\\
height = 5.7\\
is\_student = True
\end{english}

\end{codebox}

\begin{tabularx}
{\linewidth}
{
|>{\RaggedRight\arraybackslash}p{0.27\linewidth}
|>{\RaggedRight\arraybackslash}p{0.27\linewidth}
|Y|
}

\hline

\rowcolor{softBlue}

\multicolumn{1}{|c|}{\textbf{Variable}}
&
\multicolumn{1}{c|}{\textbf{Value}}
&
\multicolumn{1}{c|}{\textbf{Data-এর ধরন}}

\\

\hline

\textenglish{\ttfamily name}
&
\textenglish{\ttfamily "Rahim"}
&
Text/String
\\

\hline

\textenglish{\ttfamily age}
&
\textenglish{\ttfamily 18}
&
Integer
\\

\hline

\textenglish{\ttfamily height}
&
\textenglish{\ttfamily 5.7}
&
Floating-point
\\

\hline

\textenglish{\ttfamily is\_student}
&
\textenglish{\ttfamily True}
&
Boolean
\\

\hline

\end{tabularx}

\begin{rememberbox}

এই Data Typeগুলোর নির্দিষ্ট Python নাম, ব্যবহার ও বৈশিষ্ট্য
\textbf{5.34 Python Data Type}
অংশে বিস্তারিত আলোচনা করা হবে।

\end{rememberbox}

%================================================
% Variable Value Change
%================================================

\booksubheading{Variable-এর Value পরিবর্তন হতে পারে}

Variable-এর একটি গুরুত্বপূর্ণ ধারণা হলো—
Program চলাকালে এর Value পরিবর্তন করা যেতে পারে।

\begin{codebox}

\begin{english}
\ttfamily
x = 10\\
x = 20
\end{english}

\end{codebox}

শেষে
\textenglish{\ttfamily x}
-এর বর্তমান Value হবে—

\begin{codebox}

\begin{english}
\ttfamily
20
\end{english}

\end{codebox}

কারণ দ্বিতীয় Assignment-এ নতুন Value
আগের Value-এর পরিবর্তে ব্যবহৃত হয়েছে।

\begin{center}

\textbf{পুরোনো Value}
\(\longrightarrow\)
\textbf{নতুন Assignment}
\(\longrightarrow\)
\textbf{নতুন Value}

\end{center}

%================================================
% Assignment
%================================================

\booksubheading{Assignment-এর প্রাথমিক ধারণা}

Python-এ সাধারণভাবে
\textenglish{\ttfamily =}
চিহ্ন ব্যবহার করে Variable-এর সঙ্গে Value যুক্ত করা হয়।

উদাহরণ—

\begin{codebox}

\begin{english}
\ttfamily
number = 50
\end{english}

\end{codebox}

এখানে
\textenglish{\ttfamily =}
হলো \textbf{Assignment Operator}।

\begin{center}

\textbf{ডান পাশের Value}
\(\longrightarrow\)
\textbf{বাম পাশের Variable}

\end{center}

\begin{rememberbox}

\textbf{\color{deepOrange}গুরুত্বপূর্ণ}

Programming-এ—

\begin{center}

\textenglish{\ttfamily x = 10}

\end{center}

এর অর্থ Mathematical Equation নয়।

এখানে
\textenglish{\ttfamily 10}
Value-টি
\textenglish{\ttfamily x}
Variable-এর সঙ্গে Assign করা হচ্ছে।

\end{rememberbox}

%================================================
% Variable and Constant
%================================================

\booksubheading{Variable ও Constant-এর প্রাথমিক পার্থক্য}

Variable-এর Value পরিবর্তিত হতে পারে।

উদাহরণ—

\begin{codebox}

\begin{english}
\ttfamily
x = 10\\
x = 15
\end{english}

\end{codebox}

অন্যদিকে Constant ধারণায় কোনো Value সাধারণত
অপরিবর্তিত রাখার উদ্দেশ্যে ব্যবহার করা হয়।

\begin{rememberbox}

এখানে শুধু মূল ধারণাটি মনে রাখবে।

Python Variable-এর নিয়ম ও ব্যবহার
\textbf{5.33 Python Variable}
অংশে বিস্তারিত আলোচনা করা হবে।

\end{rememberbox}

%================================================
% Importance
%================================================

\booksubheading{Variable ও Data Type জানা কেন গুরুত্বপূর্ণ?}

Variable ও Data Type-এর ধারণা থাকলে—

\begin{compacttwocolpoints}

\item Program-এ Data সংরক্ষণ ও ব্যবহার করা সহজ হয়।

\item বিভিন্ন ধরনের Data আলাদা করে বোঝা যায়।

\item Calculation সঠিকভাবে করা যায়।

\item Input নেওয়া ও Output দেখানো সহজ হয়।

\item Condition, Loop ও Function-এ Data নিয়ে কাজ করা যায়।

\end{compacttwocolpoints}

%================================================
% Small Example
%================================================

\booksubheading{একটি ছোট উদাহরণ}

\begin{examplebox}{Program}

\begin{codebox}

\begin{english}
\ttfamily
name = "Mina"\\
age = 17\\[1mm]
print(name)\\
print(age)
\end{english}

\end{codebox}

\textbf{Output:}

\begin{codebox}

\begin{english}
\ttfamily
Mina\\
17
\end{english}

\end{codebox}

এখানে
\textenglish{\ttfamily name}
ও
\textenglish{\ttfamily age}
হলো Variable এবং তাদের সঙ্গে ভিন্ন ধরনের Value যুক্ত রয়েছে।

\end{examplebox}

%================================================
% Summary Table
%================================================

\booksubheading{এক নজরে}

\begin{tabularx}
{\linewidth}
{
|>{\bfseries\RaggedRight\arraybackslash}p{0.30\linewidth}
|Y|
}

\hline

\rowcolor{softBlue}

\multicolumn{1}{|c|}{\textbf{বিষয়}}
&
\multicolumn{1}{c|}{\textbf{মূল ধারণা}}

\\

\hline

Variable
&
Data বা Value-কে নাম দিয়ে ব্যবহার করা হয়।
\\

\hline

Value
&
Variable-এর সঙ্গে যুক্ত প্রকৃত Data।
\\

\hline

Data Type
&
Value কী ধরনের তা নির্দেশ করে।
\\

\hline

Assignment
&
Variable-এর সঙ্গে Value যুক্ত করা।
\\

\hline

Assignment Operator
&
\textenglish{\ttfamily =}
\\

\hline

Variable Value
&
প্রয়োজনে পরিবর্তন করা যায়।
\\

\hline

\end{tabularx}

%================================================
% Common Mistakes
%================================================

\booksubheading{সাধারণ ভুল}

\begin{compacttwocolpoints}

\item Variable এবং Value-কে একই বিষয় মনে করা।

\item Variable-এর নামকে Data Type মনে করা।

\item
\textenglish{\ttfamily "10"}
ও
\textenglish{\ttfamily 10}
একই ধরনের Data মনে করা।

\item
\textenglish{\ttfamily =}
-কে সবসময় Mathematical Equality মনে করা।

\item Variable-এর Value পরিবর্তন করা যায় না মনে করা।

\item Text লিখতে Quotation Mark ব্যবহার না করা।

\end{compacttwocolpoints}

%================================================
% Final Reminder
%================================================

\begin{rememberbox}

\textbf{\color{deepOrange}মনে রাখবে}

\begin{itemize}

\item
\textbf{Variable}
ব্যবহার করে Data বা Value-কে নাম দেওয়া হয়।

\item
\textbf{Value}
হলো Variable-এর সঙ্গে যুক্ত প্রকৃত Data।

\item
\textbf{Data Type}
Value কী ধরনের তা নির্দেশ করে।

\item
\textenglish{\ttfamily =}
হলো Assignment Operator।

\item
Variable-এর Value পরিবর্তিত হতে পারে।

\item
একই Program-এ বিভিন্ন ধরনের Data ব্যবহার করা যায়।

\item
দ্রুত মনে রাখবে—

\begin{center}

\textbf{Variable Name}
\(\longrightarrow\)
\textbf{Value}
\(\longrightarrow\)
\textbf{Data Type}

\end{center}

\item
উদাহরণ—

\begin{center}

\textenglish{\ttfamily age = 18}

\end{center}

এখানে—

\begin{center}

\textenglish{\ttfamily age}
\(\longrightarrow\)
\textbf{Variable}

\qquad

\textenglish{\ttfamily 18}
\(\longrightarrow\)
\textbf{Value}

\qquad

\textbf{Integer}
\(\longrightarrow\)
\textbf{Data Type}

\end{center}

\end{itemize}

\end{rememberbox}

\small

%================================================
% 5.33 Python Variable
%================================================

